\section{Visualisation Primitives}
\label{sec:theory:primitives}

The thesis of Liu~\cite{liu2022} introduces three reusable
visualisation primitives for spatial-temporal data ---
\emph{FuzzyCut}, \emph{Gridify}, and \emph{GroupSet}. Each
primitive solves a recurring class of analytical problem and
each is generic enough to be applied to several different
underlying variables. This section presents each primitive in
its own terms.

\subsection{FuzzyCut}
\label{sec:theory:fuzzycut}

Many analytical questions ask the user to categorise a
continuous variable --- a speed, a fare, a duration --- into
qualitative bins such as \emph{Slow}, \emph{Normal} and
\emph{Fast}. The choice of cut-points is rarely obvious in
advance. Fixing the cut-points produces a brittle visualisation
that misclassifies values near the boundaries, and switching
the cut-points abruptly produces distracting visual jumps as
the analyst tunes the thresholds.

FuzzyCut addresses this with \emph{fuzzy} membership functions
over the variable's domain. Each value is assigned a partial
membership in each category according to a smooth membership
curve centred on the analyst's chosen cut-point; values exactly
at the cut-point belong half to one category and half to the
neighbouring category, and the membership transitions
continuously rather than abruptly. The analyst adjusts
cut-points by dragging interactive lines on a histogram of the
variable; downstream visualisations recolour in real time, with
smooth transitions in the boundary regions rather than discrete
jumps. Liu shows that fuzzy membership reduces the visual
instability of classical discrete bins and enables analysts to
discover stable threshold settings interactively.

The technique is application-agnostic. The same FuzzyCut
mechanism can categorise traffic speeds (\emph{Congested /
Slow / Normal / Fast}), revenue per trip (\emph{High / Medium
/ Low / Loss-making}), or any other continuous variable on
which a small number of qualitative bands are useful.

\subsection{Gridify}
\label{sec:theory:gridify}

Spatial accessibility analysis asks how easily a population
can reach a given resource. The canonical question is
\emph{how far is the average resident from the nearest active
resource?} A naive answer overlays resource dots on a base map,
but this conveys density rather than accessibility: the analyst
cannot tell which areas are underserved without computing
distances explicitly.

Gridify divides the geographic region of interest into a
regular grid of configurable cell size (typically tens to
hundreds of metres on a side). Each cell is coloured by an
aggregate over the points it contains --- in the accessibility
setting, the average distance from the cell centroid to the
nearest active resource. The grid abstraction transforms a
point-density question into a region-colouring question, which
is far easier to read at a glance. A time-of-day toggle
reveals how the colouring changes across the day, exposing
resources whose availability varies with the time at which the
question is asked.

Liu demonstrates Gridify on origin--destination flows in Lyon,
but the technique applies to any spatial accessibility question
in which the analyst needs to know not where the resources are
but where they are not.

\subsection{GroupSet}
\label{sec:theory:groupset}

A common analytical question across many domains is \emph{which
elements (routes, sensors, stops, stations) behave similarly
over time?} The analyst's instinct is often to cluster the
elements by their hourly or daily time-series, but classical
clustering algorithms require the analyst to specify the number
of clusters and to fix a distance metric in advance. The result
is a single static partition that is hard to adjust as the
analyst's understanding evolves.

GroupSet replaces the single static partition with an
interactive collection of editable sets. Each element is
assigned to an initial group based on its time-series shape;
the analyst can then merge groups by dragging them together,
split a group by selecting a subset of its members and
detaching them, or relabel a group with a domain term such as
\emph{AM Peak Only} or \emph{Chronically Empty}. The technique
preserves the analyst's qualitative judgements while still
benefitting from algorithmic suggestions about which elements
group naturally together. The result is a partition that the
analyst owns, in contrast to the partition that an opaque
clustering algorithm imposes.

\subsection{Marey time--distance diagram}
\label{sec:background:marey}

A fourth primitive used in this thesis is the Marey time--distance
diagram~\cite{linzeng2014}. The diagram plots, for each trip on a
single route, a polyline on a two-dimensional grid whose
horizontal axis is wall-clock time and whose vertical axis is the
ordinal position of each stop along the route. The slope of the
polyline at any point is the bus's local speed; a horizontal
segment indicates a dwell at a stop; the vertical gap between
two consecutive polylines is the headway between two buses at
the same point on the route. Bunching, the operational pathology
in which two scheduled buses converge in time, appears in the
diagram as two polylines that move from a separated state to a
converging state along the time axis. Schedule gaps appear as
wide horizontal voids in the diagram. The Routes card of this
thesis renders a Marey diagram as its principal Detail-mode chart
(Chart~6.5 of Section~\ref{sec:proposed:dashboard}).
